\chapter{2019年江苏卷}

\section{填空题}

\begin{problem}
已知集合 \(A=\{-1,0,1,6\}\)，\(B=\{x\mid x>0,\ x\in\R\}\)，则 \(A\cap B=\)\fillinblank{}.
\end{problem}

\begin{problem}
已知复数 \((a+2\i)(1+\i)\) 的实部为 \(0\)，其中 \(\i\) 为虚数单位，则实数 \(a\) 的值是\fillinblank{}.
\end{problem}

\begin{problem}
下图是一个算法流程图，则输出的 \(S\) 的值是\fillinblank{}.

\bitmapfigure[max width=6.0cm]{img/2019/jiangsu/q3_fig1.png}
\end{problem}

\begin{problem}
函数 \(y=\sqrt{7+6x-x^2}\) 的定义域是\fillinblank{}.
\end{problem}

\begin{problem}
已知一组数据 \(6\)，\(7\)，\(8\)，\(8\)，\(9\)，\(10\)，则该组数据的方差是\fillinblank{}.
\end{problem}

\begin{problem}
从 \(3\) 名男同学和 \(2\) 名女同学中任选 \(2\) 名同学参加志愿者服务，则选出的 \(2\) 名同学中至少有 \(1\) 名女同学的概率是\fillinblank{}.
\end{problem}

\begin{problem}
在平面直角坐标系 \(xOy\) 中，若双曲线 \(x^2-\dfrac{y^2}{b^2}=1\)（\(b>0\)）经过点 \((3,4)\)，则该双曲线的渐近线方程是\fillinblank{}.
\end{problem}

\begin{problem}
已知数列 \(\{a_n\}\)（\(n\in\N^*\)）是等差数列，\(S_n\) 是其前 \(n\) 项和，若 \(a_2a_5+a_8=0\)，\(S_9=27\)，则 \(S_8\) 的值是\fillinblank{}.
\end{problem}

\begin{problem}
如图，长方体 \(ABCD-A_1B_1C_1D_1\) 的体积是 \(120\)，\(E\) 为 \(CC_1\) 的中点，则三棱锥 \(E-BCD\) 的体积是\fillinblank{}.

\bitmapfigure[max width=6.0cm]{img/2019/jiangsu/q9_fig1.png}
\end{problem}

\begin{problem}
在平面直角坐标系 \(xOy\) 中，\(P\) 是曲线 \(y=x+\dfrac4x\)（\(x>0\)）上的一个动点，则点 \(P\) 到直线 \(x+y=0\) 的距离的最小值是\fillinblank{}.
\end{problem}

\begin{problem}
在平面直角坐标系 \(xOy\) 中，点 \(A\) 在曲线 \(y=\ln x\) 上，且该曲线在点 \(A\) 处的切线经过点 \((-\e,-1)\)（\(\e\) 为自然对数的底数），则点 \(A\) 的坐标是\fillinblank{}.
\end{problem}

\begin{problem}
如图，在 \(\triangle ABC\) 中，\(D\) 是 \(BC\) 的中点，\(E\) 在边 \(AB\) 上，\(BE=2EA\)，\(AD\) 与 \(CE\) 交于点 \(O\). 若 \(\overrightarrow{AB}\cdot\overrightarrow{AC}=6\overrightarrow{AO}\cdot\overrightarrow{EC}\)，则 \(\dfrac{AB}{AC}\) 的值是\fillinblank{}.

\bitmapfigure[max width=6.0cm]{img/2019/jiangsu/q12_fig1.png}
\end{problem}

\begin{problem}
已知 \(\dfrac{\tan\alpha}{\tan\left(\alpha+\dfrac{\pi}{4}\right)}=-\dfrac23\)，则 \(\sin\left(2\alpha+\dfrac{\pi}{4}\right)\) 的值是\fillinblank{}.
\end{problem}

\begin{problem}
设 \(f(x)\)，\(g(x)\) 是定义在 \(\R\) 上的两个周期函数，\(f(x)\) 的周期为 \(4\)，\(g(x)\) 的周期为 \(2\)，且 \(f(x)\) 是奇函数. 当 \(x\in(0,2]\) 时，\(f(x)=\sqrt{1-(x-1)^2}\)，\(g(x)=\begin{cases}
k(x+2),&0<x\le 1,\\
-\dfrac12,&1<x\le 2,
\end{cases}\) 其中 \(k>0\). 若在区间 \((0,9]\) 上，关于 \(x\) 的方程 \(f(x)=g(x)\) 有 \(8\) 个不同的实数根，则 \(k\) 的取值范围是\fillinblank{}.

\bitmapfigure[max width=6.0cm]{img/2019/jiangsu/q14_fig1.png}
\end{problem}

\section{解答题}

\begin{problem}
在 \(\triangle ABC\) 中，角 \(A\)，\(B\)，\(C\) 的对边分别为 \(a\)，\(b\)，\(c\).

\begin{enumerate}
\item 若 \(a=3c\)，\(b=\sqrt2\)，\(\cos B=\dfrac23\)，求 \(c\) 的值；
\item 若 \(\dfrac{\sin A}{a}=\dfrac{\cos B}{2b}\)，求 \(\sin\left(B+\dfrac{\pi}{2}\right)\) 的值.
\end{enumerate}
\end{problem}

\begin{problem}
如图，在直三棱柱 \(ABC-A_1B_1C_1\) 中，\(D\)，\(E\) 分别为 \(BC\)，\(AC\) 的中点，\(AB=BC\).

求证：

\begin{enumerate}
\item \(A_1B_1\parallel\) 平面 \(DEC_1\)；
\item \(BE\perp C_1E\).
\end{enumerate}

\begin{center}
\bitmapfigure[max width=6.0cm]{img/2019/jiangsu/q16_fig1.png}
\end{center}
\end{problem}

\begin{problem}
如图，在平面直角坐标系 \(xOy\) 中，椭圆 \(C:\dfrac{x^2}{a^2}+\dfrac{y^2}{b^2}=1\)（\(a>b>0\)）的焦点为 \(F_1(-1,0)\)，\(F_2(1,0)\). 过 \(F_2\) 作 \(x\) 轴的垂线 \(l\)，在 \(x\) 轴的上方，\(l\) 与圆 \(F_2\)：\((x-1)^2+y^2=4a^2\) 交于点 \(A\)，与椭圆 \(C\) 交于点 \(D\). 连结 \(AF_1\) 并延长交圆 \(F_2\) 于点 \(B\)，连结 \(BF_2\) 交椭圆 \(C\) 于点 \(E\)，连结 \(DF_1\). 已知 \(DF_1=\dfrac52\).

\begin{enumerate}
\item 求椭圆 \(C\) 的标准方程；
\item 求点 \(E\) 的坐标.
\end{enumerate}

\bitmapfigure[max width=6.0cm]{img/2019/jiangsu/q17_fig1.png}
\end{problem}

\begin{problem}
如图，一个湖的边界是圆心为 \(O\) 的圆，湖的一侧有一条直线型公路 \(l\)，湖上有桥 \(AB\)（\(AB\) 是圆 \(O\) 的直径）. 规划在公路 \(l\) 上选两个点 \(P\)，\(Q\)，并修建两段直线型道路 \(PB\)，\(QA\). 规划要求：线段 \(PB\)，\(QA\) 上的所有点到点 \(O\) 的距离均不小于圆 \(O\) 的半径. 已知点 \(A\)，\(B\) 到直线 \(l\) 的距离分别为 \(AC\) 和 \(BD\)（\(C\)，\(D\) 为垂足），测得 \(AB=10\)，\(AC=6\)，\(BD=12\)（单位：百米）.

\begin{enumerate}
\item 若道路 \(PB\) 与桥 \(AB\) 垂直，求道路 \(PB\) 的长；
\item 在规划要求下，\(P\) 和 \(Q\) 中能否有一个点选在 \(D\) 处？并说明理由；
\item 对规划要求下，若道路 \(PB\) 和 \(QA\) 的长度均为 \(d\)（单位：百米），求当 \(d\) 最小时，\(P\)，\(Q\) 两点间的距离.
\end{enumerate}

\bitmapfigure[max width=6.0cm]{img/2019/jiangsu/q18_fig1.png}
\end{problem}

\begin{problem}
设函数 \(f(x)=(x-a)(x-b)(x-c)\)，\(a\)，\(b\)，\(c\in\R\)，\(f'(x)\) 为 \(f(x)\) 的导函数.

\begin{enumerate}
\item 若 \(a=b=c\)，\(f(4)=8\)，求 \(a\) 的值；
\item 若 \(a\ne b\)，\(b=c\)，且 \(f(x)\) 和 \(f'(x)\) 的零点均在集合 \(\{-3,1,3\}\) 中，求 \(f(x)\) 的极小值；
\item 若 \(a=0\)，\(0<b\le 1\)，\(c=1\)，且 \(f(x)\) 的极大值为 \(M\)，求证：\(M\le\dfrac4{27}\).
\end{enumerate}
\end{problem}

\begin{problem}
定义首项为 \(1\) 且公比为正数的等比数列为“M－数列”.

\begin{enumerate}
\item 已知等比数列 \(\{a_n\}\)（\(n\in\N^*\)）满足：\(a_2a_4=a_5\)，\(a_3-4a_2+4a_1=0\)，求证：数列 \(\{a_n\}\) 为“M－数列”；
\item 已知数列 \(\{b_n\}\) 满足 \(b_1=1\)，\(\dfrac1{S_n}=\dfrac2{b_n}-\dfrac2{b_{n+1}}\)，其中 \(S_n\) 为数列 \(\{b_n\}\) 的前 \(n\) 项和.
\begin{enumerate}
\item 求数列 \(\{b_n\}\) 的通项公式；
\item 设 \(m\) 为正整数，若存在“M－数列”\(\{c_n\}\)，对任意正整数 \(k\)，当 \(k\le m\) 时，都有 \(c_k\le b_k\le c_{k+1}\) 成立，求 \(m\) 的最大值.
\end{enumerate}
\end{enumerate}
\end{problem}

\section{附加题：解答题}

\begin{problem}[A]
已知矩阵 \(A=\begin{pmatrix}3&1\\2&2\end{pmatrix}\).

\begin{enumerate}
\item 求 \(A^2\)；
\item 求矩阵 \(A\) 的特征值.
\end{enumerate}
\end{problem}

\reuseproblemnumber
\begin{problem}[B]
在极坐标系中，已知两点 \(A\!\left(3,\dfrac{\pi}{4}\right)\)，\(B\!\left(\sqrt2,\dfrac{\pi}{2}\right)\)，直线 \(l\) 的方程为 \(\rho\sin\left(\theta+\dfrac{\pi}{4}\right)=3\).

\begin{enumerate}
\item 求 \(A\)，\(B\) 两点间的距离；
\item 求点 \(B\) 到直线 \(l\) 的距离.
\end{enumerate}
\end{problem}

\reuseproblemnumber
\begin{problem}[C]
设 \(x\in\R\)，解不等式 \(|x|+|2x-1|>2\).
\end{problem}



\begin{problem}
设 \((1+x)^n=a_0+a_1x+a_2x^2+\cdots+a_nx^n\)，其中 \(n\ge 4\)，\(n\in\N^*\). 已知 \(a_3^2=2a_2a_4\).

\begin{enumerate}
\item 求 \(n\) 的值；
\item 设 \((1+\sqrt3)^n=a+b\sqrt3\)，其中 \(a\)，\(b\in\N^*\)，求 \(a^2-3b^2\) 的值.
\end{enumerate}
\end{problem}

\begin{problem}
在平面直角坐标系 \(xOy\) 中，设点集 \(A_n=\{(0,0),(1,0),(2,0),\dots,(n,0)\}\)，\(B_n=\{(0,1),(n,1)\}\)，\(C_n=\{(0,2),(1,2),(2,2),\dots,(n,2)\}\)，\(n\in\N^*\).
令 \(M_n=A_n\cup B_n\cup C_n\). 从集合 \(M_n\) 中任取两个不同的点，用随机变量 \(X\) 表示它们之间的距离.

\begin{enumerate}
\item 当 \(n=1\) 时，求 \(X\) 的概率分布；
\item 对给定的正整数 \(n\)（\(n\ge 3\)），求概率 \(P(X\le n)\)（用 \(n\) 表示）.
\end{enumerate}
\end{problem}
